\subsection{Berechnung des Spannungsabfalls einer Diode}
% Alt: Aufgabe 3

\Aufgabe{
  \label{Aufg3}
  Gegeben ist eine Reihenschaltung aus einer Gleichspannungsquelle, einem Widerstand und einer Siliziumdiode mit einer Schleusenspannung von $U_\mathrm{S}=0,7\,\mathrm{V}$.
  Welche Spannung fällt näherungsweise über dem Widerstand $R$ ab, wenn die Versorgungsspannung $U_0=5\,\mathrm{V}$ beträgt?
  \begin{figure}[H]
    \centering
    \input{Tikz/src/FigDiodeSpannungsabfall}\alttext{Zu sehen ist eine elektrische Schaltung mit drei Bauteilen zur Bestimmung des Spannungsabfalls über einer Diode. Auf der linken Seite befindet sich die Spannungsquelle \(U_0\), deren Spannungspfeil nach unten zeigt. In Reihe zur Spannungsquelle sind ein Widerstand \(R\) und eine Diode \(D\) geschaltet. Neben der Diode befindet sich zusätzlich ein Spannungspfeil in Stromrichtung mit der Beschriftung \(U_\mathrm{D}\).}
    \label{fig:FigDiodeSpannungsabfall}
  \end{figure}
}


\Loesung{
  Da es sich um eine Siliziumdiode handelt, beträgt die Schleusenspannung und der daraus resultierende Spannungsabfall
  am Bauteil ungefähr $0,6 \text{ bis } 0,7\,\mathrm{V}$.

  In der Reihenschaltung ergibt sich der folgende Zusammenhang:
  \begin{equation*}
    U_0 = U_\mathrm{R} + U_\mathrm{D}
  \end{equation*}
  \begin{equation*}
    U_\mathrm{R} = U_0 - U_\mathrm{D}= 5\,\mathrm{V} - 0,7\,\mathrm{V} = 4,3\,\mathrm{V}
  \end{equation*}
  Siehe: Abschnitt \ref{sec:ElektrischesVerhalten}
}